% Work-MAT | 2. rocnik | pripravovana hodina 32
% Návaznost: uživatelské připravy.zip (hodiny 01–20), autorská učebnice na GitHubu.
% Realisticky.cz: didaktická inspirace v posloupnosti a typech úloh; zadání jsou vlastní.
\documentclass[11pt,a4paper]{article}
\input{../../sablony/pripravy/priprava.sty}
\nastavPripravu{2. ročník}{Logaritmy}{32}{Co je logaritmus?}
\begin{document}
\blok{Motivace a návaznost}
Z hodin o exponenciálních funkcích už umíme hledat neznámý exponent u jednoduchých mocnin. Pojmenujeme tuto operaci a navážeme na grafickou inverzi z hodin 03–06.
\poznamka{Orientační nabídka pro 45--55 minut: krátké zopakování a motivace, nový výklad, tři společné vzorové úlohy, samostatná práce. Vyber podle tempa třídy; zbývající zadání slouží jako rezerva.}
\blok{Výklad}
Zápis $\log_a b=c$ znamená právě $a^c=b$. Vždy $a>0$, $a\ne1$ a $b>0$.
\begin{definiceBox}[title={Definice - Logaritmus}] $\log_a b=c$ právě tehdy, když $a^c=b$ pro $a>0$, $a\ne1$, $b>0$.\end{definiceBox}
\blok{Společné příklady}
\begin{prikladReseniBox}[title={Příklad 1}]
Vyřeš, resp. urči: $\log_2 8=x$.
\tcblower
\[\begin{aligned}2^x=8=2^3 \\ x=3\end{aligned}\]
\end{prikladReseniBox}
\begin{prikladReseniBox}[title={Příklad 2}]
Vyřeš, resp. urči: $\log_3\frac{1}{9}=x$.
\tcblower
\[\begin{aligned}3^x=3^{-2} \\ x=-2\end{aligned}\]
\end{prikladReseniBox}
\poznamka{Každý zápis převáděj oběma směry. Základ logaritmu není hledaný exponent.}
\blok{Samostatné procvičování}
\begin{zadaniReseniBox}[title={Samostatná práce 1}]
\begin{enumerate}[label=\alph*)]
\item $\log_{2}\left(8\right)$.
\item $\log_{3}\left(9\right)$.
\item $\log_{5}\left(25\right)$.
\item $\log_{2}\left(32\right)$.
\item $\log_{4}\left(16\right)$.
\item $\log_{3}\left(27\right)$.
\end{enumerate}
\tcblower
\begin{enumerate}[label=\alph*)]
\item $3$.
\item $2$.
\item $2$.
\item $5$.
\item $2$.
\item $3$.
\end{enumerate}
\end{zadaniReseniBox}
\blok{Dobrovolné rozšíření}
\begin{zadaniReseniBox}[title={Rozšíření 1}]
\begin{enumerate}[label=\alph*)]
\item $Vytvoř vlastní úlohu tohoto typu a se spolužákem porovnej podmínky a výsledek.$.
\end{enumerate}
\tcblower
\begin{enumerate}[label=\alph*)]
\item Zkontroluj definice a argumenty..
\end{enumerate}
\end{zadaniReseniBox}
\blok{Úkol na doma}
U každého logaritmu vypiš podmínky argumentu a proveď zkoušku. U rovnic vyznač případně nepřípustné kořeny, u nerovnic odpověz intervalem.
\begin{zadaniReseniBox}[title={Domácí úkol}]
\begin{enumerate}[label=\alph*)]
\item $\log_{2}\left(\frac{1}{4}\right)$.
\item $\log_{5}\left(\frac{1}{5}\right)$.
\item $\log_{10}\left(1000\right)$.
\item $\log_{2}\left(1\right)$.
\item $\log_{3}\left(\frac{1}{9}\right)$.
\item $\log_{4}\left(64\right)$.
\item $\log_{5}\left(1\right)$.
\item $\log_{10}\left(\frac{1}{10}\right)$.
\end{enumerate}
\tcblower
\begin{enumerate}[label=\alph*)]
\item $-2$.
\item $-1$.
\item $3$.
\item $0$.
\item $-2$.
\item $3$.
\item $0$.
\item $-1$.
\end{enumerate}
\end{zadaniReseniBox}
\poznamka{Domácí úkol obsahuje 8 položek; předpokládaná domácí práce přibližně 15--30 minut. Po větších blocích jsou vědomě promíchané dříve probrané typy.}
\end{document}
